\uuid{6zsa}
\titre{Trigonométrie}
\auteur{bodin ; hanani ; mzari}
\organisation{exo7}
\licence{CC BY-NC-SA 4.0}
\datecreate{2020-01-20}
\niveau{L1}
\module{Algèbre}
\chapitre{Nombres complexes}
\sousChapitre{Forme cartésienne, forme polaire}
\difficulte{1}
\theme{Écritures algébrique et géométrique, trigonométrie}
\exo7qcmid{lille-1.nombres-complexes.006}

\contenu{
\begin{question}
\qclassification{104.01}
\qtags{motcle=trigonométrie}
Soit $\theta$ un réel.  Quelles sont les assertions vraies ?
\begin{answers}
    \bad{$\cos(2\theta)= 2\cos\theta \sin \theta$}
    \good{ $\cos(2\theta)= \cos^2\theta -\sin^2 \theta$}

    \good{$\sin(2\theta)= 2\cos\theta \sin \theta$}

    \bad{$\sin(2\theta)= \cos^2\theta -\sin^2 \theta$}
\end{answers}
\begin{explanations}
On peut appliquer la formule de Moivre, ou utiliser les formules d'addition du cosinus et du sinus.

\end{explanations}
\end{question}
}
